% !TeX root = ../main.tex
% 10. fejezet — AI-chatbot és időjárás-integráció
% Hivatkozott források: Backend/HUNting.API/Controllers/ChatController.cs,
%   Services/OpenAiService.cs, Models/Chat/ChatRequest.cs, Program.cs,
%   Frontend/HUNting/src/components/Chatbot.tsx, services/weatherService.ts,
%   pages/Weather/WeatherPage.tsx

\chapter{AI-chatbot és időjárás-integráció}
\label{chap:chatbot}

A gerinc utolsó fejezete a két külső szolgáltatásra épülő kiegészítő funkciót tárgyalja: a
mesterséges intelligenciára épülő vadászati asszisztenst és az időjárás-lekérdezést. A fejezet
egyben a funkciófejezetek egységes sablonjának határesete is: e funkciók \emph{nem használnak
saját adatbázis-entitást} -- a chatbot állapotmentes (a beszélgetés csak a kliens memóriájában
él), az időjárás pedig teljes egészében kliensoldali integráció --, ezért az adatmodell-szakasz
helyét ennek indoklása veszi át, a hangsúly pedig a külső API-integrációk mérnöki részleteire
kerül.

\section{A chatbot célja és felhasználói működése}
\label{sec:chat-cel}

A „Vadászati Asszisztens” az alkalmazás minden fő oldalán elérhető, lebegő gombból nyitható
csevegőablak, amely magyar nyelven válaszol vadászati, vadgazdálkodási és fegyverismereti
kérdésekre -- tematikusan szűkítve az 1996. évi LV. törvény (a vad védelméről, a
vadgazdálkodásról, valamint a vadászatról) és a fegyvertartási szabályozás körére, kiegészítve
vadreceptekkel és időjárási témákkal. A felhasználó kötetlen szöveges kérdést tesz fel, a válasz
formázott (címsorokat, felsorolásokat is tartalmazó) szövegként jelenik meg; a beszélgetés a
munkamenet alatt folytatólagos, az ablak bezárásáig a korábbi üzenetek is láthatók. A widget a
hitelesítési oldalakon nem jelenik meg, mert azokat külön elrendezés (\texttt{AuthLayout})
szolgálja ki, amely nem tartalmazza a komponenst.

% TODO: SC-15 képernyőkép (figures/sc-15-chatbot.png) — nyitott widget magyar válasszal; OpenAI-kulcs kell.
\begin{figure}[htbp]
	\centering
	\fbox{\parbox[c][5.5cm][c]{0.85\textwidth}{\centering SC-15: A chatbot-widget megnyitva (elkészítendő)}}
	\caption{A Vadászati Asszisztens csevegőablaka}
	\label{fig:chatbot-screen}
\end{figure}

\section{A chatbot backend-megvalósítása}
\label{sec:chat-backend}

A \texttt{ChatController} egyetlen végpontot kínál (\texttt{POST /api/Chat}), amely a
\texttt{ChatRequest} DTO-t fogadja: az új üzenetet és az eddigi beszélgetés üzenetlistáját
(szerep--tartalom párok). Üres üzenetre 400-as hibával válaszol, az integráció bármely hibáját
pedig általános, magyar nyelvű 500-as üzenetté alakítja -- a belső hibarészletek így nem
szivárognak ki a kliens felé. Az érdemi munkát az \texttt{OpenAiService} végzi, amely a
\texttt{Program.cs}-ben típusos HTTP-kliensként van regisztrálva
(\texttt{AddHttpClient<IOpenAiService, OpenAiService>}), így a \texttt{HttpClient} életciklusát a
keretrendszer gyára kezeli.

A szolgáltatás működése négy lépésre bontható. Először a konfigurációból kiolvassa az
\texttt{OpenAI:ApiKey} kulcsot; ennek hiányában naplózott kivétellel leáll -- a kulcs tehát nem a
kódban, hanem a konfigurációban (fejlesztői gépen user secrets / \texttt{appsettings}) él.
Másodszor összeállítja az üzenetlistát: az élére egy részletes magyar \emph{rendszerprompt} kerül,
amely kijelöli az asszisztens szakterületét, megköveteli a témán kívüli kérdések udvarias
elutasítását, és kifejezetten tiltja a nem létező jogszabályi adatok kitalálását. Harmadszor a
kliens által küldött előzményeket \emph{szanitálva} fűzi be: a szerepmezőt fehérlistával
ellenőrzi, és minden nem \texttt{user} vagy \texttt{assistant} értéket \texttt{user} szerepre ír
át (\ref{lst:openai-request}.~kódrészlet) -- e nélkül egy rosszindulatú kliens
\texttt{system} szerepű üzenettel felülbírálhatná a rendszerpromptot (prompt-injekció). Végül a
kérés a \texttt{gpt-4o} modellel, alacsony (0{,}2) hőmérséklettel -- azaz konzervatív,
tényközpontú beállítással -- és 1500 tokenes válaszkorláttal megy az OpenAI
chat-completions végpontjára, Bearer-fejlécben vitt kulccsal.

\begin{lstlisting}[language={[Sharp]C}, caption={Az előzmények szanitálása és a kérés összeállítása az \texttt{OpenAiService} osztályban (kivonat, megjegyzések nélkül)}, label={lst:openai-request}]
var messages = new List<object>
{
    new { role = "system", content = systemPrompt }
};

foreach (var msg in history)
{
    var role = msg.Role == "user" || msg.Role == "assistant" ? msg.Role : "user";
    messages.Add(new { role = role, content = msg.Content });
}

messages.Add(new { role = "user", content = newMessage });

var requestBody = new
{
    model = "gpt-4o",
    messages = messages,
    temperature = 0.2,
    max_tokens = 1500
};

using var request = new HttpRequestMessage(HttpMethod.Post,
    "https://api.openai.com/v1/chat/completions");
request.Headers.Authorization = new AuthenticationHeaderValue("Bearer", apiKey);
request.Content = new StringContent(JsonSerializer.Serialize(requestBody),
    Encoding.UTF8, "application/json");
\end{lstlisting}

A válaszfeldolgozás \texttt{JsonDocument}-tel navigál a
\texttt{choices[0].message.content} útvonalon; sikertelen HTTP-státusz esetén a szolgáltatás a
hibatörzset naplózza, majd általános kivételt dob, amelyet a kontroller fordít 500-as válaszra. A
beszélgetés a szerveren \emph{nem perzisztálódik}: a \texttt{ChatRequest} és a
\texttt{ChatMessage} tisztán átviteli objektum, adatbázistábla nem tartozik hozzájuk -- minden
kérés önmagában hordozza a kontextusát. A kérés teljes életútját
a~\ref{fig:chat-sequence}.~ábra mutatja.

% TODO: C-01 szekvenciadiagram (figures/c-01-chat-sequence.pdf) — Chatbot → POST /Chat →
%       OpenAiService → api.openai.com → reply; szerep-szanitálás kiemelve.
\begin{figure}[htbp]
	\centering
	\fbox{\parbox[c][6cm][c]{0.9\textwidth}{\centering C-01: A chatbot-kérés feldolgozásának szekvenciadiagramja (elkészítendő)}}
	\caption{A chatbot-kérés útja a widgettől az OpenAI API-ig és vissza}
	\label{fig:chat-sequence}
\end{figure}

\section{A chatbot frontend-megvalósítása}
\label{sec:chat-frontend}

A \texttt{Chatbot} komponens lebegő gombként ül az elrendezés bal alsó sarkában, és kattintásra
csevegőablakká nyílik. A beszélgetést egyszerű komponens-állapot tárolja: az üzenetlista egy
asszisztensi üdvözlő üzenettel indul, a beküldött kérdés azonnal bekerül a listába, a
betöltésjelző pedig a válasz megérkezéséig forog. A kérés a közös \texttt{apiClient} példányon
megy ki, az előzménylistából a nyitó üdvözlést levágva (\texttt{history: chatHistory.slice(1)}),
mert az nem a modelltől származik. Hiba esetén a widget általános hibaszöveget fűz a
beszélgetésbe, így a felhasználói élmény nem törik meg. A válaszok \texttt{react-markdown}
komponenssel renderelődnek, egyedi elemleképezéssel (címsorok, listák, kiemelések) -- a modell
által adott strukturált válaszok így olvashatóan jelennek meg. Két ergonómiai részlet teszi
teljessé: az üzenetlista minden változáskor az aljára gördül, a küldés pedig Enterrel is
kiváltható, üres bevitelnél letiltva. A beszélgetés nem kerül sem \texttt{localStorage}-ba, sem
szerverre: az oldal újratöltése új beszélgetést kezd.

\section{Időjárás-integráció}
\label{sec:weather}

Az időjárás-oldal a vadászat tervezését segíti: városnév alapján mutatja az aktuális
hőmérsékletet, hőérzetet, páratartalmat, a szöveges leírást (magyar nyelven) és a szélsebességet
-- utóbbit a szélirányt jelző, a \texttt{wind\_deg} értékkel elforgatott nyíllal. A megvalósítás
a többi funkciótól eltérő, tudatos architektúra-kompromisszum: a \texttt{weatherService} modul
\emph{közvetlenül}, a backend megkerülésével hívja az OpenWeather „current weather” végpontját
(\texttt{q}, \texttt{units=metric}, \texttt{lang=hu} paraméterekkel), saját axios-hívással, nem a
közös \texttt{apiClient}-tel -- hiszen itt sem hitelesítési fejlécre, sem a backend 401-kezelésére
nincs szükség. Az API-kulcsot a \texttt{VITE\_OPEN\_WEATHER\_API} környezeti változó adja. Az
oldal állapotkezelése szándékosan egyszerű: város, betöltés, hiba és eredmény lokális
állapotban, React Query nélkül -- egyszeri, gombnyomásra induló lekérdezésről lévén szó a
gyorsítótárazás nem adna hozzáadott értéket.

% TODO: SC-11 képernyőkép (figures/sc-11-weather.png) — találattal; OpenWeather-kulcs kell.
\begin{figure}[htbp]
	\centering
	\fbox{\parbox[c][5cm][c]{0.85\textwidth}{\centering SC-11: Az időjárás-oldal találattal (elkészítendő)}}
	\caption{Az időjárás-oldal egy lekérdezett város aktuális adataival}
	\label{fig:weather-screen}
\end{figure}

\section{Biztonsági megfontolások és korlátok}
\label{sec:chat-korlatok}

A chatbot legfontosabb védelmi vonala a szerveroldali elhelyezés: az OpenAI-kulcs soha nem jut el
a kliensig, a rendszerprompt nem módosítható kívülről, a szerep-fehérlista pedig a
prompt-injekció alapvető vektorát zárja le. Ugyanakkor a végponton \emph{nincs autorizáció és
nincs sebességkorlátozás}: anonim hívó is költséget generálhat a fizetős OpenAI-kvóta terhére --
éles üzem előtt bejelentkezéshez kötés és kérésszám-korlát (rate limiting) szükséges. Az
időjárás-integráció tükörképe ennek a kompromisszumnak: a kliensoldali hívás egyszerű és nem
terheli a backendet, cserébe az OpenWeather-kulcs a böngészőbe kerül, azaz kiolvasható -- ingyenes
kvótájú kulcsnál vállalható, éles környezetben backend-proxy javasolt. További apró, ismert
hibák: hiányzó kulcs esetén a szolgáltatás csendben \texttt{null} értékkel tér vissza, így a
felület hibajelzés nélkül marad üres, az időjárás-ikon pedig rögzített \texttt{http://} sémájú
URL-ről töltődik, amelyet a böngésző HTTPS alatt vegyes tartalomként blokkolhat. E pontok a
továbbfejlesztési irányok (\ref{chap:osszefoglalas}.~fejezet) listáját gazdagítják.
